\subsection{Correctness and Completeness}
\label{sec:formal-correctness}
We prove that the mathematical algorithm decides whether a valid assignment
exists and constructs one whenever it does.

\begin{theorem}[Complete constructive attribute synthesis]
\label{thm:complete-synthesis}
For a symbolic rewrite and metadata $m$ satisfying
Section~\ref{sec:formal-scope}, including the requirement that every input
axis have extent at least two,

\begin{align}
 \operatorname{check}(m)&\iff\exists\beta:\operatorname{Valid}(m,\beta),
 \label{eq:complete-theorem}\\
 \operatorname{check}(m)&\Rightarrow
 \operatorname{Valid}(m,\operatorname{build}(m)).
 \label{eq:sound-theorem}
\end{align}
\end{theorem}
Intuitively, \code{check} answers whether any valid RHS attribute assignment
exists for the given match metadata. If it succeeds, \code{build} produces
one; if it fails, no valid assignment exists within the defined fragment.
Thus the algorithm both finds a solution whenever possible and guarantees
that the constructed replacement is legal and equivalent to the LHS.

When a rule's \code{where} condition implies \code{check}, the constructor
supplies valid attributes for every legal match in this fragment.
Appendix~\ref{app:completeness-proof}
gives the necessity and sufficiency proofs and the extension to multiple terms.
